% Conceptual bridge. Maps each element of the classical software supply
% chain to its agent-context analogue. Pure exposition; placed after the
% Introduction so the rest of the paper can lean on the mapping. The
% right column is anchored in PrimitiveHub's own vocabulary.

\section{From Software to Agent-Context Supply Chains}
\label{sec:bridge}

The agent-context supply chain is not a metaphor borrowed loosely from
software packaging; it is the same chain with a different artifact at
each link. Classical supply-chain security hardened the path from source
code to a running binary. The agent-context path runs from a written
primitive to an action a model takes, and every defense built for the
first path has a counterpart on the second.
Table~\ref{tab:bridge} states the correspondence; the rest of the paper
fills it in.

\begin{table}[t]
\centering
\caption{Each element of the classical software supply chain has an
agent-context analogue. The right column names the artifact or control
this paper studies; parenthetical tags point to the relevant Context CI
category.}
\label{tab:bridge}
\small
\begin{tabular}{@{}p{0.46\columnwidth} p{0.46\columnwidth}@{}}
\toprule
\textbf{Software supply chain} & \textbf{Agent-context supply chain} \\
\midrule
Source code            & Context primitive \\
Build artifact         & Trust Bundle \\
SBOM                   & CBOM (context bill of materials) \\
Package registry       & Primitive registry \\
Malware signature      & Context CI detector \\
Dependency confusion   & Poisoned naming / impersonation (CCI-007) \\
Postinstall script     & Unsafe shell-install instruction (CCI-005) \\
Runtime exploit        & Tool-call / shell / credential misuse \\
Build provenance       & Trust Contract \\
Antivirus scanner      & Static semantic context scanner \\
\bottomrule
\end{tabular}
\end{table}

The correspondence is load-bearing in two directions. It lets us inherit
methodology that already works: content addressing, signed attestations,
append-only ledgers, and a labeled corpus that doubles as a regression
contract. It also marks where the analogy stops. A malware signature
matches bytes; a Context CI detector must read prose as a model would,
because the attack is carried in language, not opcodes. A postinstall
script executes on the host; an unsafe-install \emph{instruction}
executes only if a model chooses to follow it. The defense therefore
sits one step earlier than antivirus ever did: not at the moment of
execution, but at the moment a primitive is admitted to the chain.
